\section{Mechanisms}
\label{sec:mechanisms}

Sections~\ref{sec:apparent}--\ref{sec:survives} established an empirical
fact: net of volatility, observable DeFi liquidations do not amplify ETH's
downside tail to any economically significant degree. This section
rationalizes that fact through market structure. The aim is not to
re-estimate it, but to explain why an economist should not have expected
otherwise at the observable DeFi scale. The argument is a beam, not a proof.
No single channel demonstrates the null on its own, and the four are ordered
by evidentiary weight, from the load-bearing channel (size) to the modifier
(oracle). Each bounds one link in the hypothetical leverage-cycle chain,
from liquidation to directional fire sale to permanent price impact to
aggregate spiral, and the register stays descriptive throughout. The
liquidation channel is real and locally disruptive. It raises volatility, as
Section~\ref{sec:volchannel} confirms; it is only too small and too
efficiently cleared to aggregate into a directional downside loop.

\subsection{Size of the liquidation flow}
\label{sec:size}

For a liquidation flow to imprint downside amplification on ETH's
\emph{aggregate} return distribution, it must be material relative to the
turnover that forms that distribution. It is not. In the normal regime the
median daily flow is on the order of \SizeMedianShare{} of one percent of
ETH spot-plus-perpetual turnover (centralized-exchange reported), a rounding error; the mean is
\SizeShareMean\%; and even the worst day of the entire sample reaches only
\SizeShareWorstDay\% (Table~\ref{tab:sizeratio}). Against the largest reported single-day
centralized-exchange liquidation cascade on record, roughly nineteen
billion dollars on 10 October 2025 \citep{coinglass2025}, the typical
daily DeFi flow sits four to six orders of magnitude lower, and the
whole-sample cumulative total amounts to roughly \SizeTotalVsCascade\% of
that single cascade day. A flow this
many orders of magnitude below the venue where a reflexive spiral
plausibly lives (Section~\ref{sec:discussion}) cannot, at the observable
scale, be a material aggregate contributor to ETH's crash quantile. In
aggregate-market terms the hourly shock the quantile projection sees is
negligible, and that is what predicts the null.

\begin{table}[t]
\centering
\caption{Size of the DeFi liquidation flow relative to ETH turnover and to
centralized-exchange cascades.}
\label{tab:sizeratio}
% GENERATED by scripts/paper/make_tables.py — DO NOT EDIT
% fragment only: wrap in \begin{table} + caption in the manuscript
\begin{tabular}{lc}
\toprule
 & Share (\%) \\
\midrule
Mean daily DeFi liq.\ / daily ETH turnover & 0.0032 \\
Median daily DeFi liq.\ / daily ETH turnover & 0.0000 \\
Worst day (2024-08-05) / daily ETH turnover & 0.67 \\
Worst 7 days / 7-day ETH turnover & 0.10 \\
Worst day / Oct-2025 CEX-perp cascade ($\sim$\$19bn) & 1.62 \\
Whole-sample total / Oct-2025 single-day cascade & 13.3 \\
\bottomrule
\end{tabular}

\begin{minipage}{0.92\textwidth}
\vspace{2pt}
{\scriptsize \emph{Notes:} Numerator is our measured flow (the hard
number); denominators are industry turnover estimates
\citep{coingecko2024annual,coingecko2024perps} and a reported cascade
magnitude \citep{coinglass2025}, partly inflated by wash trading
\citep{gan2022} and not peer-reviewed, hence order-of-magnitude bounds. The inequality is wide enough to survive any
plausible rescaling of the denominator. Source: \texttt{size\_ratio.csv}.}
\end{minipage}
\end{table}

The inequality is wide, three to four orders of magnitude, wide enough to
survive deflating the wash-inflated denominators entirely (the numerator
is our own hard extraction; the caveat is detailed at the table). The flow
is also concentrated in a single episode (August 2024). This is why the
median, not the mean, carries the argument, and why we compare the
typical flow to a cascade rather than the outlier day to a cascade.
Agent-based work on DeFi lending finds systemic default rates below a
tenth of a percent even at ten times observed
volatility \citep{chaudhary2022},
and the ecosystem's lending base is itself non-monotone, peaking near
\$45 billion in late 2021 and falling to roughly \$15 billion by mid-2023
\citep{bertomeu2024}. ``Small'' is
not only our empirical finding but what a model of on-chain lending
predicts under stress.

\subsection{Clearing}
\label{sec:clearing}

Even a small flow could imprint directionality if it discharged as a fire
sale, with inventory sold under pressure, in size, into a falling market. The
architecture of DeFi liquidation rules this out for the modal case.
Liquidations are absorbed by well-incentivized liquidators
\citep{qin2021}; over seventy percent of liquidable positions are
liquidated immediately, and eighty-five percent of liquidable collateral
clears within two blocks \citep{perez2021}, leaving no overhang to drain
gradually into the price. And the modal liquidation is not a crash but a
drift in position health. \citet{moallemi2024} document that over
seventy percent of liquidations occur in the absence of any downward price
jump. The central mode of DeFi liquidation is non-directional; the
directional episodes are the tail, not the body. Toxic liquidation spirals
do exist as a mechanism \citep{warmuz2022},
but they are a conditional tail phenomenon (illiquid collateral, deep
undercollateralization), not the modal case, and they do not contradict
the rapid clearing of the typical liquidation.

\subsection{Absorption and price impact}
\label{sec:absorption}

The next link requires the selling pressure to move the market price and
stay there. Automated-market-maker depth and arbitrage absorb and reverse
it instead. The measured price impact of a DeFi liquidation is minuscule, on the order of $0.05\%$ on Aave and $0.01\%$ on Maker at the block level, and temporary
\citep{tian2025}. On the largest AMM, \citet{leharparlour2025} document the absence of long-lived arbitrage opportunities, so a pool-level displacement is corrected rather than propagated into a persistent move in aggregate ETH. The most serious counterweight is the
permanent price impact documented by \citet{lehar2022}, with spillover as
far as centralized venues. But their own evidence bounds its reach for our
asset, in three steps. The flagship contagion in their data runs through
an illiquid collateral token, not through deep ETH; the arbitrage they
document reverses the transient component; and their heavier tails are
explicitly two-sided, with no quantile or asymmetry test. That is a
level-impact result, not an aggregate downside-amplification one. For deep,
hourly ETH, the aggregate footprint is dominated by the symmetric volatility
channel of Sections~\ref{sec:deconstruction}--\ref{sec:survives}, not by a
directional translation of the distribution.

\subsection{Oracle decoupling}
\label{sec:oracle}

The final link would require a transient on-chain price move to feed back
into the aggregate market through the liquidation system itself. DeFi
oracles are largely centralized-exchange-anchored and smoothed: they
aggregate from external venues and are, in the words of the OECD study,
oblivious to small price changes \citep{oecd2023}. A transient
DEX deviation does not mechanically trigger a fresh liquidation wave priced
off that deviation, so there is no small-move price-oracle-liquidation
feedback loop. We place this channel last because its sign is
design-conditional. The same oracle linkage is, in other configurations,
itself a contagion channel, the broader oracle-design problem of
\citet{caldarelli2020}.
For small transient moves on a deep asset, smoothing and anchoring decouple
and attenuate; the oracle-as-cascade channel manifests on large moves of
illiquid assets, consistent with the illiquid-token flagship of
Section~\ref{sec:absorption}. The net effect is in our favor for ETH at the
hourly grain, with the acknowledgment that the reverse mechanism
exists, which is why this is a modifier and not a load-bearing channel.

The adversarial limit of this design-conditional channel is a liquidator who
does not passively receive the oracle price but actively moves it.
\citet{sadeghi2026} study profitable oracle manipulation that triggers
liquidations on constant-product markets, a form of oracle-extractable value,
and identify the transaction-fee conditions under which it becomes
economically unviable. This is the precise micro-mechanism through which a
liquidation-to-price loop could operate on chain, and a theoretical force that
already bounds it. Our result bounds it from the data side: the hourly
equivalence bound of Section~\ref{sec:boundednull} limits the asymmetric
imprint such strategies, were they prevalent on deep ETH, leave in the
aggregate tail, while staying silent on their block-level profitability, which
an hourly panel cannot resolve. The active-liquidator premise is also ours;
the flow we measure is cleared by agents who choose when to act.

\subsection{Synthesis}

The channel is real and raises volatility (Section~\ref{sec:volchannel}), but
too small relative to aggregate ETH turnover (Section~\ref{sec:size}) and too
efficiently cleared (Sections~\ref{sec:clearing}--\ref{sec:oracle}) to
aggregate into a directional downside loop in the observable on-chain channel. The null of
Section~\ref{sec:survives} does not rest on this beam; it is established
empirically and bounded by equivalence. But the beam makes it
intelligible rather than surprising, in order of evidentiary weight: size
(load-bearing, a hard number), then clearing, then absorption, then the
design-conditional oracle modifier. The account is size-contingent. It explains why the channel does not aggregate
\emph{today}, not that it never could, a conditionality
Section~\ref{sec:discussion} takes up.